\section{Second Specht branching rule for wreath products}
For $n>0$, we can embed $S_m \wr S_{n-1}$ into $S_m \wr S_n$ by mapping $(\sigma;\alpha_1,\ldots,\alpha_{n-1})$, where
$\sigma \in S_{n-1}$, $\alpha_i \in S_m$, to $ (\sigma;\alpha_1,\ldots,\alpha_{n-1},e)$, making use of the canonical embedding of
$S_{n-1}$ into $S_n$.
Hence for $\ul{\lambda} = (\lambda^1,\ldots,\lambda^r)$ an $r$-multipartition of $n$, we can consider the $k(S_m \wr S_{n-1})$-module
\[ S^{\ul{\lambda}}\bigr\downarrow^{m \wr n}_{m \wr (n-1)} \cong
T^{\ul{\lambda}}\bigr\uparrow^{m{\wr}n}_{m{\wr}|\ul{\lambda}|}\bigr\downarrow^{m \wr n}_{m \wr (n-1)}\]
obtained by restricting $S^{\ul{\lambda}}$ from $k(S_m \wr S_n)$ to $k(S_m \wr S_{n-1})$. By Mackey's Theorem we have
\begin{equation}\label{res_Sp_mwrnsub_mackey:eq}
T^{\ul{\lambda}}\bigr\uparrow^{m{\wr}n}_{m{\wr}|\ul{\lambda}|}\bigr\downarrow^{m \wr n}_{m \wr (n-1)}
\cong
\bigoplus_{u \in \mathcal{U}}\bigl(T^{\ul{\lambda}}\bigr)^u\bigl\downarrow^{(m{\wr}|\ul{\lambda}|)^u}_{(m{\wr}|\ul{\lambda}|)^u \cap m \wr (n-1)}
\bigl\uparrow^{m \wr (n-1)}_{(m{\wr}|\ul{\lambda}|)^u \cap m \wr (n-1)}
\end{equation}
with minor notational abuses as in the argument for the first branching rule, and where $\mathcal{U}$ represents a complete non-redundant
system of $(S_m \wr S_{|\ul{\lambda}|},S_m \wr S_{n-1})$-double coset representatives in $S_m \wr S_n$. We thus want to find such a set of
double coset representatives. For  $\sigma \in S_n$, let us write $\hat\sigma$ for the element $(\sigma;e,\ldots,e)$ of $S_m \wr S_n$.
Let $\sigma_1,\ldots,\sigma_N$ be a complete non-redundant system of $(S_{|\ul{\lambda}|},S_{n-1})$-double coset representatives in $S_n$. We
claim that $\hat\sigma_1,\ldots,\hat\sigma_N$ is then a complete non-redundant system of $(S_m \wr S_{|\ul{\lambda}|},S_m \wr S_{n-1})$-double
coset representatives in $S_m \wr S_n$. Indeed, if $(\theta;\alpha_1,\ldots,\alpha_n) \in S_m \wr S_n$, then we have
$\theta = \epsilon \sigma_i \delta$
for some $i \in \{1,\ldots,N\}$, $\epsilon \in S_{|\ul{\lambda}|}$ and $\delta \in S_{n-1}$, and it follows that
\[
(\theta;\alpha_1,\ldots,\alpha_n) =
\underbrace{(\epsilon;\alpha_{(1)\sigma_i},\ldots,\alpha_{(n)\sigma_i})}_{\in S_m \wr S_{|\ul{\lambda}|}}
\underbrace{(\sigma_i;e,\ldots,e)}_{=\,\hat\sigma_i}
\underbrace{(\delta;e,\ldots,e)}_{\in S_m \wr S_{n-1}}
\]
which establishes completeness. For non-redundancy, suppose that we have some $i,j$ such that
\[ \bigl(S_m \wr S_{|\ul{\lambda}|}\bigr) \hat\sigma_i  \bigl(S_m \wr S_{n-1}\bigr) =
\bigl(S_m \wr S_{|\ul{\lambda}|}\bigr) \hat\sigma_j  \bigl(S_m \wr S_{n-1}\bigr). \]
Hence $\hat\sigma_i \in \bigl(S_m \wr S_{|\ul{\lambda}|}\bigr) \hat\sigma_j  \bigl(S_m \wr S_{n-1}\bigr)$, so that
we have $\epsilon \in S_{|\ul{\lambda}|}$, $\delta \in S_{n-1}$ and elements $\alpha_i,\beta_i$ of $S_m$ such that
\[
(\sigma_i;e,\ldots,e) =
(\epsilon;\alpha_1,\ldots,\alpha_n)
(\sigma_j;e,\ldots,e)
(\delta;\beta_1,\ldots,\beta_{n-1},e)
\]
from which it follows that $\sigma_i = \epsilon\sigma_j\delta$ and hence that $i=j$. Thus we now seek such $\sigma_1,\ldots,\sigma_N$,
and to do this we shall make use of our work on tableaux.

Now recall that if $\alpha,\gamma$ are compositions of $n$, then we have defined the tableau $\tau^\alpha_\gamma$ to be the tableau of shape
$\alpha$ whose entries, read from left to right across each row in turn starting with the top row, consist of $\gamma_1$ 1's, then
$\gamma_2$ 2's, then $\gamma_3$ 3's, and so on. So for example if $n=9$, $\alpha = (8,1)$ and $\gamma = (3,1,0,2,3)$, then
\[\tau^\alpha_\gamma = \mytab{
{{1,1,1,2,4,4,5,5},{5}}
}.\]
Further, we know by Corollary \ref{tabs_give_cosets:cor} that if we have $\sigma_1,\ldots,\sigma_N \in S_n$ such that
$\tau^\alpha_\gamma\sigma_1,\ldots,\tau^\alpha_\gamma\sigma_N$ is a complete list, with no repetition, of the tableaux of shape $\alpha$ and
type $\gamma$ with weakly increasing rows, then $\sigma_1,\ldots,\sigma_N$ is in fact a complete system of $(S_\gamma,S_\alpha)$-double
coset representatives without redundancy. We now apply this in the case where $\alpha = (n-1,1)$ and $\gamma = |\ul{\lambda}|$ to obtain
our desired system of $(S_{|\ul{\lambda}|},S_{n-1})$-double coset representatives in $S_n$, noting that the subgroup $S_{n-1}$ of $S_n$ is
exactly the Young subgroup $S_{(n-1,1)}$. The following example should serve to illustrate the general argument which we shall give below.

Keep $n=9$, and suppose that $|\ul{\lambda}| = (3,1,0,2,3)$ as above. Then the possible tableaux of shape $(n-1,1)$ and type
$|\ul{\lambda}|$ with weakly increasing rows are
\[
\mytab{
{{1,1,1,2,4,4,5,5},{5}}
}\]\[
\mytab{
{{1,1,1,2,4,5,5,5},{4}}
}\]\[
\mytab{
{{1,1,1,4,4,5,5,5},{2}}
}\]\[
\mytab{
{{1,1,2,4,4,5,5,5},{1}}
}
\]
Thus, a complete non-redundant system of $(S_{|\ul{\lambda}|},S_{(n-1,1)})$-double coset representatives is
$e,(6,9,8,7),(4,9,8,7,6,5),(3,9,8,7,6,5,4)$, recalling that in our action of $S_n$ on tableaux, $\sigma \in S_n$ acts by moving the
contents of the $i$\tss{th} box to the $(i)\sigma$\tss{th} box, where the boxes of a tableau are numbered with the numbers $1,\ldots,n$
from left to right across each row, working from the top row to the bottom row.

The general case works in exactly the same way as the example. Indeed, recall that $\ul{\lambda} = (\lambda^1,\ldots,\lambda^r)$. For
$i=1,\ldots,r$ we let $b_i = |\lambda^1|+\cdots+|\lambda^i|$, so that we have a sequence $0 \leq b_1 \leq b_2 \leq \cdots \leq b_r = n$.
Then for each $i=1,\ldots,r$ such that $b_i \neq 0$ we define an element $\rho_i$ of $S_n$ by letting
\[
\rho_i = \begin{cases}
(b_i,n,n-1,\ldots,b_i+1) \quad\text{if $b_i < n$}\\
e \quad\text{if $b_i = n$}
\end{cases}
\]
(where $e$ is the identity element). By letting $i$ run through all $1,\ldots,r$ such that $|\lambda^i| > 0$, we obtain a complete list
of all the distinct $\rho_i$ without repetition. As in the above example, we see that the set of all tableaux
$\tau^{(n-1,1)}_{|\ul{\lambda}|}\rho_i$ for $i$ such that $|\lambda^i| > 0$ forms a complete list of all of the tableaux of shape $(n-1,1)$
and type $|\ul{\lambda}|$ with weakly increasing rows. Hence by Corollary \ref{tabs_give_cosets:cor} we see that the collection of all
$\rho_i$ for $i$ such that $|\lambda^i| > 0$ forms a complete non-redundant system of $(S_{|\ul{\lambda}|},S_{n-1})$-double coset
representatives in $S_n$, and hence the collection of all $\hat\rho_i$ for $i$ such that $|\lambda^i| > 0$ forms a complete non-redundant
system of $(S_m \wr S_{|\ul{\lambda}|},S_m \wr S_{n-1})$-double coset representatives in $S_m \wr S_n$.

Looking back to \eqref{res_Sp_mwrnsub_mackey:eq}, we see that we want to understand the module
\[
\bigl(T^{\ul{\lambda}}\bigr)^{\hat\rho_i}\bigl\downarrow^{(m{\wr}|\ul{\lambda}|)^{\hat\rho_i}}_{(m{\wr}|\ul{\lambda}|)^{\hat\rho_i} \cap m \wr (n-1)}
\bigl\uparrow^{m \wr (n-1)}_{(m{\wr}|\ul{\lambda}|)^{\hat\rho_i} \cap m \wr (n-1)}
\]
for $i$ such that $|\lambda^i|>0$. Our first step in doing so will be to understand the subgroup
$\bigl(S_m \wr S_{|\ul{\lambda}|}\bigr)^{\hat\rho_i}\cap\bigl(S_m \wr S_{n-1}\bigr)$ of $S_m \wr S_n$ and its action on the module
$\bigl(T^{\ul{\lambda}}\bigr)^{\hat\rho_i}$.

So choose $i$ such that $|\lambda^i|>0$.
It is easy to show directly that $\bigl(S_m \wr S_{|\ul{\lambda}|}\bigr)^{\hat\rho_i}$ is equal to
$S_m \wr \bigl(S_{|\ul{\lambda}|}\bigr)^{\rho_i}$. Thus we have
\[\bigl(S_m \wr S_{|\ul{\lambda}|}\bigr)^{\hat\rho_i}\cap\bigl(S_m \wr S_{n-1}\bigr)
=
S_m \wr \bigl(S_{|\ul{\lambda}|}\bigr)^{\rho_i}\cap\bigl(S_m \wr S_{n-1}\bigr)\]
and it is easy to show directly that $S_m \wr \bigl(S_{|\ul{\lambda}|}\bigr)^{\rho_i}\cap\bigl(S_m \wr S_{n-1}\bigr)$
is equal to the subgroup of $S_m \wr S_n$ consisting of all elements of the form
\begin{equation}\label{elt_of_conj_intersect:eq}
(\sigma; \alpha_1,\ldots,\alpha_{n-1},e)
\end{equation}
where $\sigma$ is an element of the subgroup
$\bigl(S_{|\ul{\lambda}|}\bigr)^{\rho_i}\cap S_{n-1}$ of $S_n$
and $\alpha_i \in S_m$. We thus wish to understand the subgroup $\bigl(S_{|\ul{\lambda}|}\bigr)^{\rho_i}\cap S_{n-1}$ of $S_n$.
By Proposition \ref{tab_stab:prop}, $\bigl(S_{|\ul{\lambda}|}\bigr)^{\rho_i}$ is the stabilizer (under
the action of $S_n$) of the tableau $\tau^{(n-1,1)}_{|\ul{\lambda}|}\rho_i$. It is easy to see that the tableau
$\tau^{(n-1,1)}_{|\ul{\lambda}|}\rho_i$ is the unique tableau of shape $(n-1,1)$ and type $|\ul{\lambda}|$ with weakly increasing rows which
has an $i$ in the box on the second row; such tableaux are illustrated in the above example. For any subset $\Omega$ of $\{1,\ldots,n\}$,
let us write $S(\Omega)$ to denote the subgroup of
$S_n$ consisting of all permutations which fix any number not lying in $\Omega$. We easily see that the stabilizer of the tableau
$\tau^{(n-1,1)}_{|\ul{\lambda}|}\rho_i$ is the subgroup $X^i_{|\ul{\lambda}|}$ of $S_n$, where we define (recalling that $|\lambda^i|>0$ and hence
$b_i > b_{i-1}$, where $b_0$ is taken to be 0)
\begin{multline*}
X^i_{|\ul{\lambda}|} = S\bigl(\{1,\ldots,b_1\}\bigr) \times S\bigl(\{b_1+1,\ldots,b_2\}\bigr) \times \cdots\\
\times S\bigl(\{b_{i-1}+1,\ldots,b_i-1,n\}\bigr)\times
S\bigl(\{b_i,\ldots,b_{i+1}-1\}\bigr)\times S\bigl(\{b_{i+1},\ldots,b_{i+2}-1\}\bigr)\times\\
\cdots \times S\bigl(\{b_{r-1},\ldots,b_r-1 =n-1\}\bigr)
\end{multline*}
(note that here we are using the $\times$ symbol to denote an \textit{internal} direct product of subgroups, and that if $b_i = b_{i+1}$
then $\{b_i,\ldots,b_{i+1}-1\}$ represents the empty set, and that if $b_i = b_{i-1}+1$ then $\{b_{i-1}+1,\ldots,b_i-1,n\} = \{n\}$), and
hence $\bigl(S_{|\ul{\lambda}|}\bigr)^{\rho_i} = X^i_{|\ul{\lambda}|}$.
We now introduce a small piece of notation.
Indeed, if $\gamma = (\gamma_1,\ldots,\gamma_r)$ is a composition of $n$, and $i \in \{1,\ldots,r\}$ such that $\gamma_i > 0$, then we
write $[\gamma]_i$ for the composition $(\gamma_1,\ldots,\gamma_{i-1},\gamma_i -1, \gamma_{i+1},\gamma_r)$ of $n-1$.
We see that $X^i_{|\ul{\lambda}|}\cap S_{n-1}$ is the subgroup
\begin{multline*}
S\bigl(\{1,\ldots,b_1\}\bigr) \times S\bigl(\{b_1+1,\ldots,b_2\}\bigr) \times \cdots\\
\times S\bigl(\{b_{i-1}+1,\ldots,b_i-1\}\bigr)\times S\bigl(\{b_i,\ldots,b_{i+1}-1\}\bigr)\times S\bigl(\{b_{i+1},\ldots,b_{i+2}-1\}\bigr)
\times\\
\cdots \times S\bigl(\{b_{r-1},\ldots,b_r-1 =n-1\}\bigr)
\end{multline*}
of $S_n$, and under our embedding of $S_{n-1}$ into $S_n$
this is exactly the subgroup $S_{[|\ul{\lambda}|]_i}$ of $S_{n-1}$.
Hence, recalling that we are viewing $S_m \wr S_{n-1}$ as a subgroup
of $S_m \wr S_n$ via the embedding $(\sigma;\alpha_1,\ldots,\alpha_{n-1}) \longmapsto (\sigma;\alpha_1,\ldots,\alpha_{n-1},e)$, we see that
the subgroup $\bigl(S_m \wr S_{|\ul{\lambda}|}\bigr)^{\hat\rho_i}\cap\bigl(S_m \wr S_{n-1}\bigr)$ of $S_m \wr S_n$ is equal to the subgroup
$S_m \wr S_{[|\ul{\lambda}|]_i}$ of the subgroup $S_m \wr S_{n-1}$ of $S_m \wr S_n$.

We now turn our attention to the action of
$\bigl(S_m \wr S_{|\ul{\lambda}|}\bigr)^{\hat\rho_i}\cap\bigl(S_m \wr S_{n-1}\bigr)$
on the
$k\bigl(S_m \wr S_{|\ul{\lambda}|}\bigr)^{\hat\rho_i}$-module $\left(T^{\ul{\lambda}}\right)^{\hat\rho_i}$.
We know by the definition of conjugate modules that $\left(T^{\ul{\lambda}}\right)^{\hat\rho_i}$ is the module formed by equipping
$T^{\ul{\lambda}}$ with the $k\bigl(S_m \wr S_{|\ul{\lambda}|}\bigr)^{\hat\rho_i}$-action $\ast$ given for $x \in T^{\ul{\lambda}}$ and
$y \in \bigl(S_m \wr S_{|\ul{\lambda}|}\bigr)^{\hat\rho_i}$ by $x \ast y = x (\hat\rho_iy\hat\rho_i^{-1})$ (where the action on the
right-hand side is the action of $S_m \wr S_{|\ul{\lambda}|}$ on $T^{\ul{\lambda}}$, noting that $\hat\rho_iy\hat\rho_i^{-1}$ does indeed
lie in $S_m \wr S_{|\ul{\lambda}|}$). Thus to calculate the action of an element
\[(\sigma;\alpha_1,\ldots,\alpha_{n-1},e) \in \bigl(S_m \wr S_{|\ul{\lambda}|}\bigr)^{\hat\rho_i}\cap\bigl(S_m \wr S_{n-1}\bigr)\]
on the module $\left(T^{\ul{\lambda}}\right)^{\hat\rho_i}$, we need to calculate
$\hat\rho_i(\sigma;\alpha_1,\ldots,\alpha_{n-1},e)\hat\rho_i^{-1}$. We have
\begin{align*}
\hat\rho_i(\sigma;\alpha_1,\ldots,\alpha_{n-1},e)\hat\rho_i^{-1}
&= (\rho_i;e,\ldots,e)(\sigma;\alpha_1,\ldots,\alpha_{n-1},e)(\rho_i^{-1};e,\ldots,e)\\
&= (\rho_i\sigma\rho_i^{-1};\alpha_{(1)\rho_i},\ldots,\alpha_{(n)\rho_i})\qquad \text{(taking $\alpha_n = e$)}\\
&= (\rho_i\sigma\rho_i^{-1};\alpha_1,\alpha_2,\ldots,\alpha_{b_i-1},e,\alpha_{b_i},\alpha_{b_i+1},\\
&\hspace{13em}\ldots,\alpha_{n-2},\alpha_{n-1}).
\end{align*}
But by our description \eqref{elt_of_conj_intersect:eq} of the elements of
$\bigl(S_m \wr S_{|\ul{\lambda}|}\bigr)^{\hat\rho_i}\cap\bigl(S_m \wr S_{n-1}\bigr)$, we see that
$\sigma \in \bigl(S_{|\ul{\lambda}|}\bigr)^{\rho_i}\cap S_{n-1}$, which implies that
$\rho_i\sigma\rho_i^{-1} \in S_{|\ul{\lambda}|}\cap \bigl(S_{n-1}\bigr)^{\rho^{-1}_i}$. By direct calculation, any element of
$\bigl(S_{n-1}\bigr)^{\rho^{-1}_i}$ fixes $b_i$, and hence we see that $\rho_i\sigma\rho_i^{-1}$ is an element of
$S_{|\ul{\lambda}|}$ which fixes $b_i$.
Now we know that the subgroup $S_{|\ul{\lambda}|}$ of $S_n$ has an internal direct product factorisation
\begin{multline*}
S\bigl(\{1,\ldots,b_1\}\bigr) \times S\bigl(\{b_1+1,\ldots,b_2\}\bigr) \times \cdots\\
\cdots\times S\bigl(\{b_{i-1}+1,\ldots,b_i\}\bigr)\times S\bigl(\{b_i+1,\ldots,b_{i+1}\}\bigr)\times\cdots\\
\cdots \times S\bigl(\{b_{r-1}+1,\ldots,b_r = n\}\bigr).
\end{multline*}
Thus any element $\pi$ of $S_{|\ul{\lambda}|}$ has a unique factorisation $\pi = \theta_1\cdots\theta_r$ where
$\theta_j \in S\bigl(\{b_{j-1}+1,\ldots,b_j\}\bigr)$ (with $b_0$ taken to be 0). We thus see that $\rho_i\sigma\rho_i^{-1}$ has such a
factorisation $\rho_i\sigma\rho_i^{-1} = \theta_1\cdots\theta_r$, where $\theta_i$ fixes $b_i$. 
Thus we see that our element
$(\sigma;\alpha_1,\ldots,\alpha_{n-1},e)$ of $\bigl(S_m \wr S_{|\ul{\lambda}|}\bigr)^{\hat\rho_i}\cap\bigl(S_m \wr S_{n-1}\bigr)$ acts on the
module $\bigl(T^{\ul{\lambda}}\bigr)^{\hat\rho_i}$ as the element
\[
\bigl(\theta_1\cdots\theta_r;\alpha_1,\alpha_2,\ldots,\alpha_{b_i-1},e,\alpha_{b_i},\alpha_{b_i+1},\ldots,\alpha_{n-2},\alpha_{n-1}\bigr)
\]
of $S_m \wr S_{|\ul{\lambda}|}$ acts on $T^{\ul{\lambda}}$ (recalling that $\bigl(T^{\ul{\lambda}}\bigr)^{\hat\rho_i}$ and $T^{\ul{\lambda}}$ are equal
as $k$-vector spaces). But we know that $\bigl(S_m \wr S_{|\ul{\lambda}|}\bigr)^{\hat\rho_i}\cap\bigl(S_m \wr S_{n-1}\bigr)$ is equal to the
subgroup $S_m \wr S_{[|\ul{\lambda}|]_i}$ of the subgroup $S_m \wr S_{n-1}$ of $S_m \wr S_n$, and we now see that if we identify
$S_m \wr S_{[|\ul{\lambda}|]_i}$ with
\begin{multline*}
(S_m \wr S_{|\lambda^1|})\times(S_m \wr S_{|\lambda^2|})\times\cdots\\
\cdots\times
(S_m \wr S_{|\lambda^{i-1}|})\times(S_m \wr S_{|\lambda^i|-1})\times(S_m \wr S_{|\lambda^{i+1}|})\times\cdots\times
(S_m \wr S_{|\lambda^r|})
\end{multline*}
in the canonical way, then by the definition of the $k(S_m \wr S_{|\ul{\lambda}|})$-module $T^{\ul{\lambda}}$, the
$k(S_m \wr S_{[|\ul{\lambda}|]_i})$-module
\[
\bigl(T^{\ul{\lambda}}\bigr)^{\hat\rho_i}\bigl\downarrow^{(m{\wr}|\ul{\lambda}|)^{\hat\rho_i}}_{(m{\wr}|\ul{\lambda}|)^{\hat\rho_i} \cap m \wr (n-1)}
=
\bigl(T^{\ul{\lambda}}\bigr)^{\hat\rho_i}\bigl\downarrow^{(m{\wr}|\ul{\lambda}|)^{\hat\rho_i}}_{m \wr [|\ul{\lambda}|]_i}
\]
is isomorphic to
\begin{multline}\label{T_ul_lambd_isom_in_place_res:eq}
\left(\bigl(S^{\mu^1}\bigr)^{\widetilde\boxtimes |\lambda^1|}\oslash S^{\lambda^1}\right)\boxtimes\cdots\boxtimes
\left(\bigl(S^{\mu^i}\bigr)^{\widetilde\boxtimes |\lambda^i|}\oslash S^{\lambda^i}\right)\Bigl\downarrow^{m \wr |\lambda^i|}_{m \wr (|\lambda^i|-1)}
\boxtimes\cdots\\
\cdots\boxtimes
\left(\bigl(S^{\mu^r}\bigr)^{\widetilde\boxtimes |\lambda^r|}\oslash S^{\lambda^r}\right).
\end{multline}
Thus, we want to investigate the $k(S_m \wr S_{|\lambda^i|-1})$-module
\[\left(\bigl(S^{\mu^i}\bigr)^{\widetilde\boxtimes |\lambda^i|}\oslash S^{\lambda^i}\right)\Bigl\downarrow^{m \wr |\lambda^i|}_{m \wr (|\lambda^i|-1)}.\]
Now the restriction operation $\Bigl\downarrow^{m \wr |\lambda^i|}_{m \wr (|\lambda^i|-1)}$ may be expressed as
\[\Bigl\downarrow^{m \wr |\lambda^i|}_{m \wr (|\lambda^i|-1,1)}\Bigl\downarrow^{m \wr (|\lambda^i|-1,1)}_{m \wr (|\lambda^i|-1)},\]
where, we recall, $m \wr (|\lambda^i|-1,1)$ represents the subgroup $S_m \wr S_{(|\lambda^i|-1,1)}$ of $S_m \wr S_{|\lambda^i|}$ consisting of all
elements of the form $(\sigma;\alpha_1,\ldots,\alpha_n)$ for $\alpha_i \in S_m$ and $\sigma \in S_{(|\lambda^i|-1,1)}$, while
$m \wr (|\lambda^i|-1)$ represents the subgroup $S_m \wr S_{(|\lambda^i|-1)}$ of $S_m \wr S_{|\lambda^i|}$ consisting of all elements of the form
$(\sigma;\alpha_1,\ldots,\alpha_{n-1},e)$ for $\alpha_i \in S_m$ and $\sigma \in S_{(|\lambda^i|-1,1)}$. Now we have by
Proposition \ref{wreath_ind_res:prop} that
\[
\left(\bigl(S^{\mu^i}\bigr)^{\widetilde\boxtimes |\lambda^i|}\oslash S^{\lambda^i}\right)\Bigl\downarrow^{m \wr |\lambda^i|}_{m \wr (|\lambda^i|-1,1)}
=
\bigl(S^{\mu^i}\bigr)^{\widetilde\boxtimes |\lambda^i|}\Bigl\downarrow^{m \wr |\lambda^i|}_{m \wr (|\lambda^i|-1,1)}\:\oslash\:
S^{\lambda^i}\Bigl\downarrow^{|\lambda^i|}_{(|\lambda^i|-1,1)}.
\]
Upon further restriction to $S_m \wr S_{(|\lambda^i|-1)}$, we see that this is isomorphic to the direct sum of $\mathrm{dim}_k(S^{\mu^i})$
copies of the module
\begin{equation}\label{one_factor_S_down_to_limin1:eq}
\bigl(S^{\mu^i}\bigr)^{\widetilde\boxtimes |\lambda^i|-1}\,\oslash\,S^{\lambda^i}\!\Bigl\downarrow^{|\lambda^i|}_{|\lambda^i|-1}.
\end{equation}
It now follows by Theorem \ref{jam_branch_rule_sn:thm} and the fact that $- \oslash -$ preserves submodule series (see above)
that, if we let $\delta^{t} \triangleright \delta^{t-1} \triangleright \cdots \triangleright \delta^{1}$ be all the (distinct) partitions
of $|\lambda^i|-1$ which can be obtained by removing a single box from $\lambda^i$, then \eqref{one_factor_S_down_to_limin1:eq} has a
submodule series whose quotients are, from top to bottom,
\[
\bigl(S^{\mu^i}\bigr)^{\widetilde\boxtimes |\lambda^i|-1}\,\oslash\,S^{\delta^{t}},\ldots,
\bigl(S^{\mu^i}\bigr)^{\widetilde\boxtimes |\lambda^i|-1}\,\oslash\,S^{\delta^{1}}.
\]
Using \eqref{T_ul_lambd_isom_in_place_res:eq}, it now follows that the $k(S_m \wr S_{[|\ul{\lambda}|]_i})$-module
\[\bigl(T^{\ul{\lambda}}\bigr)^{\hat\rho_i}\bigl\downarrow^{(m{\wr}|\ul{\lambda}|)^{\hat\rho_i}}_{(m{\wr}|\ul{\lambda}|)^{\hat\rho_i} \cap m \wr (n-1)}
=
\bigl(T^{\ul{\lambda}}\bigr)^{\hat\rho_i}\bigl\downarrow^{(m{\wr}|\ul{\lambda}|)^{\hat\rho_i}}_{m \wr [|\ul{\lambda}|]_i}\]
is the direct sum of $\mathrm{dim}_k(S^{\mu^i})$ copies of a module $V_i$ which has a submodule series whose quotients are, from top to
bottom $T^{\ul{\delta}^t}, \ldots, T^{\ul{\delta}^1}$, where $\ul{\delta}^l$ is the multipartition of $n-1$ obtained by replacing $\lambda^i$
with $\delta^l$ in $\ul{\lambda}$. By exactness of the functor $\bigl\uparrow^{m \wr (n-1)}_{m \wr [|\ul{\lambda}|]_i}$, it now follows that the
$k(S_m \wr S_{n-1})$-module
\[
\bigl(T^{\ul{\lambda}}\bigr)^{\hat\rho_i}\bigl\downarrow^{(m{\wr}|\ul{\lambda}|)^{\hat\rho_i}}_{(m{\wr}|\ul{\lambda}|)^{\hat\rho_i} \cap m \wr (n-1)}
\bigl\uparrow^{m \wr (n-1)}_{(m{\wr}|\ul{\lambda}|)^{\hat\rho_i} \cap m \wr (n-1)}
=
\bigl(T^{\ul{\lambda}}\bigr)^{\hat\rho_i}\bigl\downarrow^{(m{\wr}|\ul{\lambda}|)^{\hat\rho_i}}_{m \wr [|\ul{\lambda}|]_i}
\bigl\uparrow^{m \wr (n-1)}_{m \wr [|\ul{\lambda}|]_i}
\]
is the direct sum of $\mathrm{dim}_k(S^{\mu^i})$ copies of the module $X_i = V_i\bigl\uparrow^{m \wr (n-1)}_{m \wr [|\ul{\lambda}|]_i}$, and that
$X_i$ has a submodule series whose factors are, from top to bottom, $S^{\ul{\delta}^t},\ldots,S^{\ul{\delta}^1}$.
Referring back to the decomposition \eqref{res_Sp_mwrnsub_mackey:eq}, we now see that we have proved the following result, which is our
desired Specht branching rule.
